\chapter{Les tickets, carnet de recherche}

Un projet de recherche accumule des questions ouvertes : une référence à vérifier, une page
illisible, un lemme qui manque, une lettre à écrire. Dans un carnet, elles se
perdent entre deux calculs. Les services d'hébergement offrent pour cela des
\emph{tickets} (\emph{issues}), conçus pour suivre des bogues et qui servent très bien à
suivre des questions.

\section{Une question par ticket}

\begin{regle}
Un ticket pose une seule question, assez précise pour qu'on sache quand elle est réglée.
\end{regle}

« Vérifier la bibliographie » ne se ferme jamais. « Vérifier que le théorème 2.4 de
l'article cité dit bien ce qu'on lui fait dire, et donner la page » se ferme le jour où
c'est fait. Le titre énonce la question ; le corps donne le contexte, les liens, et ce qui
a déjà été essayé.

\section{Listes de cases}

Quand une question se découpe en étapes, le corps du ticket porte une liste de cases à
cocher. L'hébergeur en affiche l'avancement, et chaque case cochée dit ce qui est fait sans
qu'on ait à relire le fil.

\begin{terminal}
- [x] lire l'exposé I, chercher les énoncés précis
- [x] recherche bibliographique, énoncé par énoncé
- [ ] formaliser ce qui n'est pas dans la littérature
- [ ] proposer le résultat à la bibliothèque commune
\end{terminal}

\section{Relier les tickets}

Un ticket en appelle d'autres. Écrire \texttt{\#89} dans le texte d'un ticket crée un lien
vers le ticket 89, et l'hébergeur montre le lien des deux côtés. Un commit dont le message
contient \texttt{\#89} apparaît dans le fil du ticket : la question et le travail qui la
règle se retrouvent au même endroit.

\begin{exemple}
Une suite de trois tickets peut ainsi garder la trace d'une recherche : le premier constate
qu'une notation change de sens au milieu d'un chapitre ; le deuxième recense les passages
touchés ; le troisième, ouvert plus tard, demande à un spécialiste de trancher un cas
douteux. Chacun renvoie au précédent.
\end{exemple}

\section{Fermer}

\begin{regle}
On ferme un ticket quand sa question est réglée, et on dit comment dans un dernier
message.
\end{regle}

Le dernier message (« réglé par le commit \texttt{4233568} : la référence est la
proposition 5.1 de l'article cité ») rend le ticket utile à qui le retrouvera. Un ticket
qu'on abandonne se ferme aussi, avec la raison de l'abandon : un abandon expliqué est une
information.

\section{Les tickets comme mémoire publique}

Sur un dépôt public, les tickets sont lisibles par tous. Ils montrent comment le travail a
été fait, ce qui a été vérifié, ce qui reste douteux. C'est un avantage pour la confiance
qu'on peut accorder au résultat. C'est aussi une contrainte : on n'y écrit rien qu'on ne
publierait pas, ni citation longue d'un texte protégé, ni information personnelle
(chapitre~\ref{ch:droits}).
